\refstepcounter{section}
\section*{Lecture 1: Introduction}
\addcontentsline{toc}{section}{Lecture 1: Introduction}
We will give a general introduction to what glassy systems are. These are characterised by some effective disorder, unlike crystals. Moreover, these systems are characterised by `slow dynamics' and are thus said to be \textit{frozen} as some specific degrees of freedom are not evolving. Sometimes the dynamics is so slow that the phase space is not effectively explored. However, small vibrational motions may be present in these systems. \\[0.2cm]
Glassy systems do not possess any definite symmetry or structure, they are \textit{amorphous}. Hence, the usual solid state techniques cannot be applied nicely in these systems!
There are different types of glasses like spin glasses, Fermi glass, Bose glass etc. We will first focus on structural glass.
